\documentclass[conference]{IEEEtran}
\IEEEoverridecommandlockouts
\usepackage{cite}
\usepackage{amsmath,amssymb,amsfonts}
\usepackage{algorithmic}
\usepackage{graphicx}
\usepackage{textcomp}
\usepackage{xcolor}
\usepackage{hyperref}
\usepackage{booktabs}
\usepackage{listings}

\def\BibTeX{{\rm B\kern-.05em{\sc i\kern-.025em b}\kern-.08em
    T\kern-.1667em\lower.7ex\hbox{E}\kern-.125emX}}

\lstset{
  basicstyle=\ttfamily\footnotesize,
  breaklines=true,
  frame=single,
  backgroundcolor=\color{gray!10}
}

\begin{document}

\title{Context Sovereignty in Autonomous Coding Agents: A 3-Layer Pruning Pipeline for Long-Running Sessions}

\author{\IEEEauthorblockN{1\textsuperscript{st} Masih Moafi}
\IEEEauthorblockA{\textit{Independent Researcher}\\
masihmoafi12@gmail.com}
}

\maketitle

\begin{abstract}
Large Language Models (LLMs) used as autonomous coding agents face a critical limitation: context-window exhaustion during long-running sessions. As agents explore codebases, execute shell commands, and read files, transient exploration data rapidly fills the context window, leaving little room for durable project knowledge and active task goals. This leads to degraded performance, high inference costs, and frequent session restarts. We introduce \textit{Elpis}, a terminal environment that enforces Context Sovereignty by explicitly decoupling the active model working set from the durable session transcript. Elpis implements a novel 3-layer pruning pipeline---comprising pre-model shell output filtering, deterministic safety truncation, and post-turn meaning-aware pruning---that distills completed tasks into compact evidence pointers. Alongside a user-facing Context Ledger for explicit admission control and a 2-stage memory architecture for curating cross-session facts, Elpis achieves a 93\% free context ratio in long-running tasks compared to 73\% in baseline unmanaged coding agents. Our architecture demonstrates that shifting context management from the model's internal attention mechanism to an explicit, inspectable external pipeline significantly improves agent longevity and reliability.
\end{abstract}

\begin{IEEEkeywords}
Autonomous Agents, Large Language Models, Context Management, Software Engineering, Prompt Engineering
\end{IEEEkeywords}

\section{Introduction}
The adoption of Large Language Models (LLMs) for autonomous software engineering tasks has grown exponentially \cite{bubeck2023sparks}. Tools and environments allow agents to act autonomously by executing shell commands, reading files, and writing code \cite{tian2023chatgpt}. However, these agents struggle with long-running sessions due to the fixed, albeit growing, size of their context windows.

When an agent is asked to debug a complex issue or implement a multi-file feature, it generates a significant amount of transient data: failed grep searches, verbose compiler outputs, and iterative file reads. In standard architectures, this transcript is appended to the context window verbatim. As the session progresses, the context fills with dead ends and raw logs \cite{liu2024lost}. The critical information---the user's original goal, durable architectural rules, and the current task state---becomes buried, leading to degraded model reasoning and increased token costs.

To address this, we present \textit{Elpis}, a provider-neutral terminal user interface (TUI) and environment for coding agents. Elpis operates on the principle of \textit{Context Sovereignty}: the agent's working context is treated as a strictly budgeted working set, not a raw chat transcript. By explicitly managing what enters and leaves the context window, Elpis ensures that the agent retains focus over long-running sessions.

Our core contributions are:
\begin{itemize}
    \item A \textbf{3-Layer Pruning Pipeline} that systematically removes transient exploration data from the working context, replacing it with compact evidence pointers while preserving full transcripts on disk.
    \item A \textbf{Context Ledger} that provides human developers with real-time, explicit control over exactly which files, rules, and checkpoints are admitted to the LLM's context.
    \item A \textbf{Two-Stage Curated Memory} system that separates ephemeral turn data from durable project facts, ensuring long-term continuity across model switches and session restarts.
    \item Empirical evaluation showing Elpis maintains 93\% free context at the end of a complex workflow, compared to 73\% for unmanaged baselines.
\end{itemize}

\section{Background and Related Work}
Current autonomous coding agents largely rely on raw transcript accumulation. Systems like AutoGPT \cite{autogpt2023} and generic CLI wrappers pass the entire session history back to the model on each turn. While models with 1M+ token context windows (e.g., Gemini 1.5 Pro) exist, empirical studies show that LLMs suffer from the ``Lost in the Middle'' phenomenon, where reasoning degrades when critical information is buried in massive contexts \cite{liu2024lost}.

Furthermore, passing large contexts incurs high latency and financial costs. Approaches to mitigate this include vector database retrieval (RAG) \cite{lewis2020retrieval} and model-driven summarization. However, RAG often struggles to maintain the sequential narrative of an active debugging session, and naive summarization frequently drops critical file paths and state variables necessary for code execution. Elpis differentiates itself by treating context management as an explicit pipeline of deterministic and meaning-aware pruning steps external to the LLM itself.

\section{System Architecture}

Elpis is designed as an assimilation environment. The underlying language model (e.g., GPT-4, Claude 3.5, Gemini) acts as the reasoning engine, but Elpis owns the context, memory, permissions, and tool execution layer.

\subsection{Context Lifetimes}
We categorize all information into three distinct lifetimes:
\begin{enumerate}
    \item \textbf{Durable:} Global rules (\texttt{AGENTS.md}), active goals (\texttt{GOAL.md}), and cross-session curated memory (\texttt{MEMORY.md}). This survives restarts and compaction.
    \item \textbf{Task:} Active thread state, decisions, changed paths, and the Executive Summary checkpoint (\texttt{ES.md}).
    \item \textbf{Turn:} Transient exploration such as file reads, directory listings, and raw command outputs. These expire immediately after the active turn.
\end{enumerate}

\subsection{The 3-Layer Pruning Pipeline}
The core of Elpis's context management is its 3-layer pipeline, executed around the model's inference loop.

\subsubsection{Layer 1: Pre-Model Shell Filtering}
Verbose shell commands (e.g., \texttt{rg}, \texttt{find}) produce massive outputs. Layer 1 uses an optional deterministic filter (RTK) to compact this output using pattern matching before the agent ever sees it, often reducing raw output by 72-97\%.

\subsubsection{Layer 2: Deterministic Safety Cap}
A hard deterministic upper bound truncates exceptionally large tool output blobs to protect context limits. Truncation notices are injected at the boundaries so the model is aware of the missing data.

\subsubsection{Layer 3: Post-Turn Meaning-Aware Pruning}
After the model completes a turn, Elpis evaluates the exploration history. Useful results are distilled into a compact conclusion and an exact evidence pointer (e.g., a rollout ID or log path). Dead ends and raw terminal outputs are evicted from the next working context entirely. The full, unaltered transcript remains on disk for human auditing.

\subsection{Context Ledger and Admission Control}
Elpis introduces the Context Ledger, a real-time TUI panel that lists every admitted portable context source alongside its exact byte size and percentage of the model's context window. Developers can dynamically toggle admission for specific rule files or goals. This ensures that context inclusion is a deliberate action, rather than an invisible side-effect of conversation length.

\subsection{Session Continuity and Curated Memory}
To survive thread compaction and model switches, Elpis supports \textit{Lean Continuation}. Instead of replaying transcripts, a fresh thread is anchored using \texttt{GOAL.md} (the objective) and \texttt{ES.md} (the latest checkpoint containing final results and changed files).

Additionally, a two-stage memory pipeline runs asynchronously. Stage 1 extracts candidate facts from completed rollouts. Phase 2 (Consolidation) evaluates these candidates against recall frequency and query diversity, promoting stable project facts to a durable \texttt{MEMORY.md}. Crucially, when an agent recalls a memory concerning workspace state, Elpis treats it as a hypothesis and re-verifies it against the live filesystem.

\section{Evaluation}

\subsection{Experimental Setup}
To evaluate the efficacy of the 3-Layer Pruning Pipeline, we compared Elpis against a standard, unmanaged coding agent baseline (Codex-derived architecture without Elpis pruning). Both systems were given the same complex software engineering task and the same initial prompt, utilizing the same underlying LLM infrastructure.

% TODO (Author): Insert specific methodology details here (e.g., what was the exact task? How many files were touched? What model was used?)

\subsection{Context Efficiency}
We measured the percentage of the context window that remained free at the end of the completed workflow. A higher free context percentage indicates that the system successfully discarded transient exploration data, leaving room for future tasks and reducing inference costs.

As shown in Table \ref{tab:results}, Elpis maintained 93\% free context at the conclusion of the workflow. In contrast, the baseline system retained only 73\% free context, having accumulated the entire transcript of searches, reads, and dead ends.

\begin{table}[htbp]
\centering
\caption{Free Context Remaining at Workflow Conclusion}
\label{tab:results}
\begin{tabular}{@{}lc@{}}
\toprule
\textbf{System} & \textbf{Free Context at End} \\ \midrule
Baseline (Unmanaged Codex) & 73\% \\
\textbf{Elpis (3-Layer Pruning)} & \textbf{93\%} \\ \bottomrule
\end{tabular}
\end{table}

\subsection{Qualitative Observations}
% TODO (Author): Expand this section with observations on latency, cost savings, or qualitative improvements in agent reasoning due to a cleaner context.
Beyond the quantitative reduction in context size, we observed that the Elpis agent was less prone to hallucinating obsolete file paths that were present early in the session but subsequently changed. By relying on the \texttt{ES.md} checkpoint, the agent maintained a highly accurate representation of the current workspace state.

\section{Conclusion}
The accumulation of transient exploration data is a major bottleneck for autonomous coding agents. Elpis demonstrates that Context Sovereignty---achieved via a 3-layer pruning pipeline, explicit admission control, and a curated two-stage memory architecture---drastically reduces context bloat. By maintaining 93\% free context after complex tasks, Elpis enables long-running sessions, reduces inference costs, and ensures that the agent remains focused on the user's goals rather than its own raw history. Future work will explore dynamic thresholding for automated tier-switching between exact resume and lean continuation.

\bibliographystyle{IEEEtran}
\bibliography{references}

\end{document}
